\documentclass[10pt,a4paper]{amsart}
\usepackage[margin=19mm]{geometry}
\usepackage{amsmath,amssymb,mathtools}
\usepackage[scr=rsfs,cal=euler]{mathalfa}
\usepackage{xfrac}
\usepackage[unicode,hidelinks,bookmarks=false]{hyperref}
\newcommand{\ie}{i.e.\ }
\newcommand{\cf}{cf.\ }
\newcommand{\resp}{respectively\ }
\newcommand{\Subsection}{\subsection}
\providecommand{\nobreakdash}{\nobreak\hbox{-}}
\setlength{\emergencystretch}{2em}
\multlinegap0pt
\usepackage{bm}
\usepackage{bbm}
\usepackage{dsfont}
\usepackage{xspace}
\usepackage{cancel}
\usepackage{mathtools}
\usepackage{amsthm}
\makeatletter
\@ifundefined{theorem}{\newtheorem{theorem}{Theorem}}{}
\@ifundefined{definition}{\newtheorem{definition}[theorem]{Definition}}{}
\@ifundefined{proposition}{\newtheorem{proposition}[theorem]{Proposition}}{}
\makeatother
\title[From Flows to Maps: Sampling Laws for Attractor Intensity an]{From Flows to Maps: Sampling Laws for Attractor Intensity and Bounded-Noise Escape}
\author{Jiguang Yu, Louis Shuo Wang}
\date{Source: arXiv:2608.02933v1 (CC BY 4.0)}
\begin{document}
\maketitle

\noindent\emph{Source and licence.} From Flows to Maps: Sampling Laws for Attractor Intensity and Bounded-Noise Escape. Version arXiv:2608.02933v1 (CC BY 4.0), distributed under the Creative Commons Attribution 4.0 International licence, as recorded in the arXiv licence metadata for this exact version and shown on the abstract page https://arxiv.org/abs/2608.02933v1. Full rights notice: licenses/arxiv-2608.02933-CC-BY-4.0.txt.\par\smallskip

\noindent\textit{Editorial scope.} Counted unit: the Proposition whose source label is \texttt{prop:scalar-one-sided} in the article (source0.tex lines 1373--1402, source label \texttt{prop:scalar-one-sided}), delivered together with the article's own complete original proof of it (source0.tex lines 1404--1520, one \texttt{proof} environment of 117 lines). No mathematics is written, repaired or re-derived: statement and proof are the article's own text line for line. Required context retained uncounted: eq:autonomous-system (lines 331--334); def:continuous-intensity (lines 409--423). Every definition, display and label of those lines is retained unchanged. Omitted from this package and not redistributed: 1--330, 335--408, 424--1372, 1403--1403, 1521--3573. Nothing inside a retained excerpt is omitted or altered: every retained body is delivered exactly as the article's lines are written, so no omission receipt of any kind accompanies this package. Citations: the retained excerpts carry no citation command at all, so no bibliography entry of the article is transcribed and no reference is redistributed. Reference adaptation: none was needed and none was made; every reference inside the delivered unit resolves inside it, so no unresolved reference or citation is produced. Printed slips kept literal: no printed word, symbol, hypothesis or number of the counted statement or of its proof was altered. Layout and class: the article's own class is \texttt{sn-jnl} with packages including geometry, graphicx, threeparttable, multirow, amssymb, amsmath, amsfonts, amsthm, mathrsfs, mathtools, appendix, xcolor, textcomp, manyfoot; the wrapper is the lane's pinned \texttt{amsart} editorial wrapper at 10pt on A4, which restates the theorem-like environments the retained text needs and adds the article's own macro definitions that the retained text needs (none), with extra wrapper packages (bm, bbm, dsfont, xspace, cancel, mathtools). The delivered numbering is the unit's own. Assets: no retained span contains a figure environment, a graphics-inclusion command, a PDF-page inclusion, a table or a TikZ picture; no image file is copied into this package, the package declares no asset dependency and no image file is redistributed with it. Licence: version arXiv:2608.02933v1 of this article is distributed under the Creative Commons Attribution 4.0 International licence, as recorded in the arXiv licence metadata for this exact version and shown on the abstract page https://arxiv.org/abs/2608.02933v1. A full rights notice is delivered at licenses/arxiv-2608.02933-CC-BY-4.0.txt. Characters: every retained excerpt is pure ASCII, so the delivered engine, XeLaTeX with its default Latin Modern text font, reports no missing character.
\begin{equation}
\label{eq:autonomous-system}
    \dot x=f(x),\qquad x\in\mathbb R^d,
\end{equation}

\begin{definition}[Continuous-time intensity]
\label{def:continuous-intensity}
Let $A$ be an attractor of \eqref{eq:autonomous-system}. Its continuous-time intensity is
\begin{equation}
\label{eq:continuous-intensity}
    \mu(A)
    :=
    \sup\left\{
    r\geq0:
    \overline{\mathcal P_r(A)}
    \subset K\Subset\mathcal D(A)
    \text{ for some compact }K
    \right\}.
\end{equation}
\end{definition}

\begin{proposition}[One-sided sampling bounds for scalar systems]
\label{prop:scalar-one-sided}
Let $f:\mathbb R\to\mathbb R$ be locally Lipschitz with forward-complete
flow, let $b<a$, and assume:
\begin{enumerate}
\item $f(b)=f(a)=0$, $f>0$ on $(b,a)$, and $f\leq0$ on $[a,\infty)$;
\item $A=\{a\}$ is an attractor with basin
$\mathcal D(A)=(b,\infty)$;
\item for every $r\geq0$ there exists $x^+(r)>a$ such that
$f(x)\leq-r$ for all $x\geq x^+(r)$.
\end{enumerate}
Set $I:=[b,a]$, let $L\geq0$ be a Lipschitz constant for $f$ on $I$,
and write
    $\mu:=\max_{x\in I}f(x)$.
Then, for every $h>0$,
\begin{align}
\label{eq:scalar-one-sided-continuous}
    \mu(A)&=\mu,\\
\label{eq:scalar-one-sided-discrete}
    \mu_h(A)&=\max_{x\in I}\bigl(F_h(x)-x\bigr),\\
\label{eq:scalar-one-sided-bounds}
    \mu\,\frac{1-e^{-Lh}}{Lh}
    \;\leq\;
    \frac{\mu_h(A)}{h}
    \;&\leq\;
    \mu,
\end{align}
where the prefactor in \eqref{eq:scalar-one-sided-bounds} is
interpreted as $1$ when $L=0$.
\end{proposition}

\begin{proof}
Throughout, we use that controlled solutions are continuous in time and
that $F_h$ is nondecreasing, being the time-$h$ map of a scalar flow.

\emph{Step 1: barriers.}
Let $r\geq0$ and let $c\in(b,a)$ satisfy $f(c)>r$. No solution of
$\dot x=f(x)+u(t)$ with $\|u\|_\infty<r$ crosses $c$ from right to
left. Indeed, suppose $x(t_0)\geq c$ and $x(t_1)<c$ for some
$t_1>t_0$, and let
$\tau:=\sup\{t\in[t_0,t_1]:x(t)\geq c\}$, so that $x(\tau)=c$ and
$x<c$ on $(\tau,t_1]$. By continuity there is $\sigma\in(\tau,t_1]$
with $f(x(t))\geq\tfrac12(f(c)+r)$ for $t\in[\tau,\sigma]$, whence
\[
    x(t)-c
    =
    \int_\tau^t\bigl(f(x(s))+u(s)\bigr)\,ds
    \geq
    (t-\tau)\,\tfrac12\bigl(f(c)-r\bigr)
    >0
    \qquad\text{on }(\tau,\sigma],
\]
contradicting the definition of $\tau$. The mirrored argument, using
hypothesis (iii), shows that no such solution crosses $x^+(r)$ from
left to right.

\emph{Step 2: continuous intensity.}
Let $r<\mu$ and let $x_*\in(b,a)$ maximize $f$ on $I$, so
$f(x_*)=\mu>r$. By Step 1, every admissible controlled trajectory
starting at $a$ remains in $[x_*,x^+(r)]$, a compact subset of
$(b,\infty)=\mathcal D(A)$. Hence every $r<\mu$ is admissible in
Definition~\ref{def:continuous-intensity}, and $\mu(A)\geq\mu$.
Conversely, let $r>\mu$ and choose $\mu<r'<r$. Under the constant
control $u\equiv-r'$, one has $\dot x=f(x)-r'\leq\mu-r'<0$ on $I$, so
the solution starting at $a$ reaches $b\notin\mathcal D(A)$ in finite
time. Thus $r$ is not admissible, and $\mu(A)\leq\mu$.

\emph{Step 3: discrete intensity.}
Set $\varepsilon_h:=\max_{x\in I}(F_h(x)-x)$. Since $F_h(x)>x$ on
$(b,a)$ and $F_h-\operatorname{id}$ vanishes at $b$ and $a$, the
maximum is positive and is attained at some $x_h\in(b,a)$.

First, $\varepsilon_h$ is admissible. Consider a perturbed orbit from
$a$ with kicks $\|e_k\|<\varepsilon_h$. If $x\geq x_h$, monotonicity
gives
    $F_h(x)+e
    \geq
    F_h(x_h)-\|e\|
    >
    x_h+\varepsilon_h-\varepsilon_h
    =
    x_h$,
so the orbit remains to the right of $x_h>b$. For confinement on the
right, pick $r$ with $rh\geq\varepsilon_h$ and set
$M:=x^+(r)+\varepsilon_h$. For $x\geq x^+(r)$, the unperturbed
trajectory decreases at rate at least $r$ while it remains in
$[x^+(r),\infty)$, and it cannot re-enter that set once it leaves it
(Step 1 with $u\equiv0$); hence
$F_h(x)\leq\max\bigl(x-rh,\;x^+(r)\bigr)\leq\max(x-\varepsilon_h,\,
M-\varepsilon_h)$. By monotonicity, for every $x\leq M$,
    $F_h(x)+e
    \leq
    F_h(M)+\varepsilon_h
    \leq
    M$,
so $(-\infty,M]$ is forward invariant for the perturbed dynamics.
Every perturbed orbit from $a$ therefore remains in the compact set
$[x_h,M]\subset(b,\infty)$, and $\mu_h(A)\geq\varepsilon_h$.

Conversely, let $\varepsilon>\varepsilon_h$ and choose
$\varepsilon_h<\varepsilon'<\varepsilon$. Apply the constant kicks
$e_k\equiv-\varepsilon'$. For $x\in(b,a]$,
\[
    F_h(x)-\varepsilon'
    \leq
    x+\varepsilon_h-\varepsilon'
    =
    x-(\varepsilon'-\varepsilon_h),
\]
while for $x\geq a$ one has $F_h(x)\leq x$, because the flow is
nonincreasing on $[a,\infty)$ and cannot cross the equilibrium $a$.
Starting at $a$, the perturbed orbit therefore decreases by at least
$\varepsilon'-\varepsilon_h$ per step while it remains in $(b,a]$, and
after finitely many steps it produces a point $\leq b$, which lies
outside $\mathcal D(A)$. Hence $\varepsilon$ is not admissible, and
$\mu_h(A)\leq\varepsilon_h$. This proves
\eqref{eq:scalar-one-sided-discrete}.

\emph{Step 4: the bounds.}
For $x\in(b,a)$, the unperturbed trajectory $\phi^s(x)$ increases
toward $a$ and hence remains in $I$. Therefore
    $F_h(x)-x
    =
    \int_0^h f\bigl(\phi^s(x)\bigr)\,ds
    \leq
    h\mu$,
and taking the maximum over $I$ gives $\mu_h(A)\leq h\mu$, the upper
bound in \eqref{eq:scalar-one-sided-bounds}.

For the lower bound, set
\[
    g(s):=f\bigl(\phi^s(x_*)\bigr)\geq0,
    \qquad
    G(s):=\int_0^s g(\sigma)\,d\sigma
    =
    \phi^s(x_*)-x_*.
\]
Since $\phi^s(x_*)\in I$, the Lipschitz bound on $I$ gives
    $g(s)
    \geq
    f(x_*)-L\bigl(\phi^s(x_*)-x_*\bigr)
    =
    \mu-LG(s)$,
so $(e^{Ls}G)'=e^{Ls}(G'+LG)\geq\mu e^{Ls}$, and integration yields
    $G(h)\geq\mu\,\frac{1-e^{-Lh}}{L}$.
Since $\mu_h(A)\geq F_h(x_*)-x_*=G(h)$ by
\eqref{eq:scalar-one-sided-discrete}, the lower bound follows.
\end{proof}

\end{document}
